%% Standalone supplementary material for RQ4 (external-validity human panel).
%% Compiles to its own PDF; NOT part of the 8-page main_v2.tex submission.
%% Covers, in full, what main_v2.tex Sec. "External-validity set" only summarizes:
%%   1. full data-generation chain, 2. operational label definitions,
%%   3. exact annotator materials, 4. no-majority handling,
%%   5. system-output-to-action mapping (Table 4), 6. exact metric formulas,
%%   plus the ethics paragraph in full.
\documentclass[11pt]{article}
\usepackage[margin=1in]{geometry}
\usepackage{booktabs,amsmath,hyperref,enumitem}
\title{Supplementary Material: RQ4 External-Validity Protocol\\
\large VerifiQuant human-labeled messy-query panel}
\author{Anonymous Author(s)}
\date{}

\begin{document}
\maketitle
\noindent This document is the full protocol behind Sec.~4.4/5.4 (``RQ4: Does an independent
human panel corroborate the boundary?'') of the main paper. The main paper states the
headline design and results; this supplement is the reproducible, auditable record.

\section{Full data-generation chain}
The 82-item messy-query set is derived from FinanceReasoning test questions in three stages,
none of which involves VerifiQuant, its contracts, or its gate taxonomy.

\begin{enumerate}[leftmargin=*]
\item \textbf{Cross-model authoring.} Each item starts from a FinanceReasoning source question
(\texttt{source\_id}, \texttt{source\_question}, \texttt{source\_gold}) used only as topical
grounding, not copied. An independent LLM author, \textbf{GPT-5.2 (OpenAI)}, rewrites it into a
single natural-sounding user message. This author model is deliberately \emph{different from
the solver stack} (Gemini 2.5 Flash, used by every evaluated system) so that the messiness of a
question is not correlated with what any solver finds easy or hard, and is never shown
VerifiQuant's contracts, gates, or the M/N/F/E/I/C taxonomy at any point.
\item \textbf{Intent-stratified generation.} Each source question is rewritten under one of five
personas (a hurried sell-side analyst, a retail investor unsure of terminology, a non-finance
operations manager, a precise quant, a headline reader sanity-checking a claim) crossed with one
of four \emph{natural user-behavior intents}, not VerifiQuant operators:
  \begin{itemize}[leftmargin=*]
  \item \emph{clean}: every needed figure is stated, one specific number is requested.
  \item \emph{missing}: the user pastes a near-complete excerpt (2--5 labeled numbers) but
  naturally omits exactly one indispensable value, with no placeholder, flag, or announcement of
  the omission.
  \item \emph{ambiguous}: all values are present, but the wording is loose about which exact
  interpretation is wanted (e.g., which metric, which period, which basis), such that more than
  one number is defensible. Plain one-period year-over-year change is explicitly excluded as
  \emph{not} ambiguous.
  \item \emph{judgment}: the user wants an opinion, a decision, or a forecast, or asks for
  something requiring a model or data not obtainable from the text.
  \end{itemize}
  These four intents are a sampling stratification, not a label; the intent used to generate an
  item is never shown to annotators, and the majority human label often disagrees with it (see
  Sec.~2 of the audit trail, \texttt{docs/2026-08-03\_rq4\_messy\_external\_validity.md}, for the
  intent-vs-majority cross-tabulation).
\item \textbf{Deterministic curation for the ``missing'' stratum.} 18 of the 82 items (the
\emph{missing}-intent items) were produced or audited by
\texttt{scripts/curate\_messy\_missing\_items.py}, an offline, deterministic fallback that
replaces or validates the LLM author's output against two hard constraints: at least two visible
labeled numeric inputs, and exactly one documented indispensable omission
(\texttt{omitted\_input}, \texttt{why\_required} in \texttt{messy\_pilot\_provenance.json}). The
omission audit is written only to provenance and is never shown on the labeling sheet or to any
solver.
\end{enumerate}
Outputs: \texttt{messy\_pilot.jsonl} (\texttt{id, question, context, persona, intended,
source\_id, author\_model}), \texttt{messy\_pilot\_provenance.json} (adds
\texttt{source\_gold} and, for missing items, the omission audit),
\texttt{messy\_pilot\_labeling.\{csv,md\}} (the blank labeling sheet, Sec.~3).
No solver is run at generation time; all four action labels are assigned by humans (Sec.~2--3),
after generation is frozen.

\section{Operational definitions of the four labels}
Annotators assign exactly one of four labels per item, following a top-down decision procedure
(first ``yes'' wins):
\begin{enumerate}[leftmargin=*]
\item Is the \emph{ask itself} not ``compute one number''? (advice/opinion, a forecast, or it
needs a method/data/model unobtainable from the text) $\to$ \textbf{O -- OUT-OF-SCOPE}.
\item Is a required numeric \emph{value} missing from the text? $\to$ \textbf{R --
REQUEST-INFO}.
\item Are all values present, but would an \emph{interpretation} choice (percent vs.\ decimal,
begin vs.\ end period, which metric, 360 vs.\ 365, gross vs.\ net, per-share vs.\ total) change
the number? $\to$ \textbf{C -- CLARIFY}.
\item Otherwise: complete, unambiguous, one determinate number $\to$ \textbf{A -- ANSWER}.
\end{enumerate}
Two contrasts annotators are explicitly trained on:
\begin{itemize}[leftmargin=*]
\item \textbf{R vs.\ C}: R means a \emph{number} is missing (ask for the value); C means all
numbers are present but \emph{which interpretation} is unclear (ask which one).
\item \textbf{C vs.\ O}, the one-question test: \emph{``If I could ask exactly one clarifying
question about interpretation, could I then compute one exact number from the data given?''}
Yes $\to$ C. No $\to$ O.
\end{itemize}
Golden rules given to annotators (abridged; full text and worked examples in
\texttt{verifiquant/data/runs/messy\_query\_pilot/LABELING\_GUIDELINE.md} and its Chinese
translation, both shipped in the repository):
\begin{enumerate}[leftmargin=*]
\item Judge from the question and context only, as the assistant would; ignore any system's
design or output (no system was shown).
\item Do not infer an unstated number from nearby figures. A partial excerpt that omits one
required value is still \textbf{R}, even if other numbers are present.
\item User confusion is not task ambiguity. ``I'm not sure how to compute it,'' with all data
present and a standard method, is \textbf{A}.
\item The label is what the assistant should \emph{do}, not whether the question's own claim is
correct. A well-posed ``does this check out?'' verification is \textbf{A} (compute the true
value and flag any discrepancy).
\item Scope is about the question, not the annotator's own skill. An advanced but fully
determined computation (duration VaR, bond convexity, option Greeks) is \textbf{A}, not
\textbf{O}.
\item One label per item, with a one-line rationale; a genuinely undecidable item may be flagged
with \texttt{?} rather than forced, and is routed to the disagreement pool (Sec.~4).
\end{enumerate}

\section{What annotators saw, and how they labeled}
Each annotator worked from an identical blank labeling sheet
(\texttt{messy\_pilot\_labeling.csv}, columns: \texttt{id, question, context, label, reason}),
containing \emph{only} the rewritten question and its context string for all 82 items. No
annotator was shown: the source FinanceReasoning question or gold answer, the generating
persona, the generating intent stratum (clean/missing/ambiguous/judgment), any system's output,
or any other annotator's labels. Annotators filled the \texttt{label} column with one of
A/R/C/O (optionally suffixed \texttt{?}) and a one-line \texttt{reason}, working independently
in their own copy of the sheet, with explicit instructions not to compare with one another until
each had finished all 82 items. Five annotators (finance graduates of the authors' institution)
completed the task; per-rater coverage was 82/82 for four annotators and 81/82 for the fifth, who
correctly left one item (\texttt{mq-58}) as a genuine \texttt{?} rather than a forced guess (that
item still reached a $4/5$ majority; Sec.~4).

\section{No-majority handling}
An item's ground-truth label is its strict majority vote ($\geq\lceil n/2\rceil{+}1$ of the
raters who gave a decision on that item, i.e.\ $\geq 3$ of 5 in the typical case; a rater's blank
or lone \texttt{?} does not count as a vote and is excluded from that item's denominator, not
imputed). Items that do \emph{not} reach a majority are written to a separate
\texttt{disagreement\_pool.csv} for manual adjudication (by the authors, disclosed as such) before
any system is scored, and are excluded from Fleiss' $\kappa$'s modal rater-count group if their
rater count differs from the majority group. In the frozen run reported in the paper, all 82
items reached a majority on the first pass (disagreement pool size $=0$), so no adjudication was
required; the mechanism above is retained here for reproducibility and because it was exercised
during protocol design on pilot batches.
This entire step (\texttt{scripts/aggregate\_messy\_labels.py}) runs before, and independently
of, any system output: \texttt{ground\_truth.jsonl} is frozen first, and only then are the four
systems scored against it.

\section{System-output-to-action mapping (full Table 4 protocol)}
Table~\ref{tab:appendix-actionmap} gives the exact, code-level mapping from each system's raw
output to one of the four action labels $\{A,R,C,O\}$, used to compute Match and Human-FNAR
(Sec.~6). $\mathrm{emit}=1$ means the system committed a numeric answer (as opposed to
abstaining, asking for a value, or refusing).

\begin{table}[h]
\centering
\caption{Raw-output $\to$ action mapping, by system.}
\label{tab:appendix-actionmap}
\small
\begin{tabular}{@{}lll@{}}
\toprule
System & Raw signal & Action rule \\
\midrule
CoT, R-PAL (plain) & committed numeric answer present? &
  $\mathrm{emit}{=}1\Rightarrow$A; else action unresolved (no R/C/O signal exposed) \\
R-PAL$^{+}$ & \texttt{assessment.decision}/\texttt{.diagnostic} &
  (execute,READY)$\to$A; (warn,I\_SOFT)$\to$A; (refuse,N)$\to$O; \\
 & & (clarify,F)$\to$R; (clarify,M)$\to$O; (clarify,I\_HARD)$\to$C; \\
 & & unparseable assessment (2/82 items) $\to$ unresolved \\
VerifiQuant & \texttt{status}/\texttt{diagnostic\_type} &
  status=success$\to$A (incl.\ I\_SOFT-warned successes); \\
 & & diagnostic\_type=F$\to$R (takes priority over a generic \texttt{status=error}); \\
 & & diagnostic\_type$\in\{$N,M,E$\}\to$O \\
\bottomrule
\end{tabular}
\end{table}

\paragraph{Two conventions made explicit, and why.}
(i) \emph{I\_SOFT counts as emit.} A VerifiQuant I\_SOFT exit still executes and returns a number
under a stated warning; the paper's own funnel semantics (``I\_SOFT proceeds with a warning and C
computes'') make this a committed answer, so it is scored as $\mathrm{emit}{=}1$, action A, for
both R-PAL$^{+}$ and VerifiQuant.
(ii) \emph{M is scored as O, not C, even though every M-refusal in this run carries an explicit
clarification question.} This follows the main paper's stated funnel-outcome convention
(\textsc{out-of-scope}$\leftrightarrow$N/M refusal) rather than the item's literal behavior, to
keep the paper's own stated mapping internally consistent. This is conservative: an
\texttt{M$\to$C} scoring was computed for comparison and raises VerifiQuant's Match from 31.7\%
to 34.1\%; it is reported here, not substituted into the main table, precisely to avoid
special-casing the convention after seeing which one scores higher.

\paragraph{4-way vs.\ binary Match.} CoT and plain R-PAL never expose \emph{why} they did not
answer (no R/C/O-level signal), only whether they answered. Scoring them on a 4-way action match
would silently discard credit for correct abstention on non-A items (an abstained item's
action is \texttt{None}, which cannot equal any human label). Match for these two systems is
therefore \emph{binary} answer-vs-abstain agreement with the human majority action
($1[\mathrm{emit}_S(i) = 1[\mathrm{GT}(i){=}\text{A}]]$). R-PAL$^{+}$ and VerifiQuant expose a
real R/C/O-level decision and are scored on the full four-way action match.

\section{Exact metric formulas}
Let $\mathrm{GT}(i)\in\{A,R,C,O\}$ be the human majority label for item $i$, over the
$n{=}82$ items. Let $\mathrm{RCO} = \{i : \mathrm{GT}(i)\neq A\}$, the set of items a legitimate
answer is a plain number is inappropriate for ($|\mathrm{RCO}|{=}46$).

\paragraph{Human-FNAR (headline metric; lower is safer).}
\[
\mathrm{HumanFNAR}(S) = \frac{\left|\{\, i \in \mathrm{RCO} : \mathrm{emit}_S(i){=}1 \,\}\right|}{|\mathrm{RCO}|}
\]
The share of items a human panel says should \emph{not} get a plain number, on which system $S$
emits one anyway.

\paragraph{Match.}
\[
\mathrm{Match}(S) =
\begin{cases}
\dfrac{\left|\{\, i : \mathrm{action}_S(i) = \mathrm{GT}(i) \,\}\right|}{n} & \text{if } S \text{ exposes a real R/C/O signal (R-PAL}^{+}\text{, VerifiQuant)} \\[1.2em]
\dfrac{\left|\{\, i : \mathrm{emit}_S(i) = \mathbb{1}[\mathrm{GT}(i){=}A] \,\}\right|}{n} & \text{otherwise (CoT, R-PAL)}
\end{cases}
\]
An item with an unresolved action (e.g.\ an unparseable assessment) counts as a non-match; it is
never imputed or excluded from the denominator.

\paragraph{Selective accuracy (Sel.\,Acc).}
\[
\mathrm{SelAcc}(S) = \frac{\left|\{\, i : \mathrm{emit}_S(i){=}1 \,\wedge\, \mathrm{NumMatch}(v_S(i), g(i)) \,\}\right|}{\left|\{\, i : \mathrm{emit}_S(i){=}1 \,\}\right|}
\]
where $v_S(i)$ is $S$'s emitted numeric value, $g(i)$ is the item's gold value (joined from
\texttt{messy\_pilot\_provenance.json}; not present for R/C/O items, so an answered R/C/O item
can never contribute a correct match), and $\mathrm{NumMatch}$ reuses the pipeline's existing
numeric comparator (\texttt{run\_error\_classification\_pipeline.\_answer\_match}, no new
tolerance introduced): exact/near-exact match at $1\%$ relative or $0.01$ absolute tolerance; an
integer-rounding mode when the question explicitly asks to round to the nearest integer; and a
$\times100$ equivalence for a decimal-vs.-percentage-point output-scale mismatch. Raw system
outputs that are not directly parseable as a float (2 of 82 items, both from R-PAL$^{+}$: one
percent-formatted string, one dict-shaped free-text answer) are recovered via a best-effort
numeric parse (strip \texttt{\$}/\texttt{\%}/commas, else take the first numeric token in the
string) rather than discarded, so that a correct-but-oddly-formatted answer is not penalized.

\section{Ethics, in full}
The 82 annotated questions are derivative of FinanceReasoning, a public benchmark, rewritten by
an LLM author into synthetic natural-language phrasing; no item contains personal or
identifiable information about any real individual. The annotation task collected exactly two
pieces of information per item per annotator: a categorical label (A/R/C/O, optionally
\texttt{?}) and a one-line rationale about \emph{the question}, not about the annotator. No
demographic, contact, or other personal data about the five annotators was collected at any
point; they are referred to throughout the paper, this supplement, and the released data only as
\emph{Annotator 1--5}, with no names in any committed file (annotator CSVs are stored as
\texttt{A1.csv}--\texttt{A5.csv}). Participation was voluntary; annotators are acknowledged as
contributors, and the task involved no compensation drawn from any grant or funding
instrument. Given (i) the absence of any research funding or contract governing this specific
task, (ii) the exclusively non-personal, publicly-derived nature of the annotated material, and
(iii) the absence of any collected data about the annotators themselves, this task does not meet
the definition of human-subjects research under the authors' institution's research-ethics
review framework; this determination was confirmed in consultation with the authors' academic
advisor prior to release. No IRB/REC approval number is claimed, and none was sought, because
none was determined to be required.

\end{document}
